\section{本讲总框架：从性能现象到可执行优化}

这节课的核心不是“学一个新 API”，而是学会把一个 GPU 程序拆成可以测、可以看、可以改的几个层次。你可以把整讲内容压成五个问题：
\begin{enumerate}
    \item \textbf{定义}：什么叫 benchmark、profile、kernel、fusion、PTX？
    \item \textbf{直觉}：为什么很多操作慢，不是算得慢，而是搬运太多？
    \item \textbf{机制}：GPU、CUDA、Triton、torch.compile 在底层分别做了什么？
    \item \textbf{应用}：GeLU、Softmax、MLP、MatMul 分别怎么优化？
    \item \textbf{局限}：什么时候理论不够，必须回到真实硬件上测？
\end{enumerate}

\begin{figure}[H]
\centering
\begin{tikzpicture}[node distance=0.95cm, every node/.style={font=\small, align=center}]
\node[draw, rounded corners, fill=blue!10, minimum width=3.1cm, minimum height=0.8cm] (a) {benchmark};
\node[draw, rounded corners, fill=blue!18, right=1.2cm of a, minimum width=3.1cm, minimum height=0.8cm] (b) {profile};
\node[draw, rounded corners, fill=blue!26, right=1.2cm of b, minimum width=3.1cm, minimum height=0.8cm] (c) {kernel};
\node[draw, rounded corners, fill=blue!34, right=1.2cm of c, minimum width=3.1cm, minimum height=0.8cm] (d) {fuse};
\draw[-{Latex[length=3mm]}, thick] (a) -- (b);
\draw[-{Latex[length=3mm]}, thick] (b) -- (c);
\draw[-{Latex[length=3mm]}, thick] (c) -- (d);
\end{tikzpicture}
\caption{本讲的工作流：先测时间，再找瓶颈，再读内核，再决定要不要融合}
\end{figure}

\begin{knowledgebox}{这讲最重要的结论}
不要先凭感觉优化。先 benchmark，再 profile，再决定是该用 PyTorch、Triton、CUDA，还是干脆交给 \texttt{torch.compile}。
\end{knowledgebox}

\subsection{定义}

Benchmark 是端到端计时。Profile 是看时间花在哪。Kernel 是 GPU 上的执行单元。Fusion 是把多个算子合并成一个内核，减少中间结果读写。PTX 是 Triton/CUDA 最终落到 GPU 之前的中间汇编视图。

\begin{importantbox}{如果只记一个定义}
benchmark 回答“多快”，profile 回答“为什么快/慢”，kernel 回答“机器到底执行了什么”。
\end{importantbox}

\subsection{直觉}

GPU 性能的核心直觉很简单：计算本身往往没那么贵，贵的是把数据从 DRAM 反复搬来搬去。只要你把多个操作分散成多个 kernel，中间写回全局内存的次数就会迅速上升。

\begin{warningbox}{最常见的直觉错误}
很多人会盯着 FLOPs 看，但在本讲里，真正决定性能的常常是 memory traffic 和 kernel launch 的开销，而不是纯计算量。
\end{warningbox}

\subsection{机制}

在机制层面，本讲会依次把工具链串起来：
\begin{itemize}
    \item PyTorch 负责表达高层算子。
    \item CUDA 让你直接写线程级内核。
    \item Triton 让你在 block 级别写 Python 内核。
    \item \texttt{torch.compile} 把朴素 Python 代码自动编译成更优的底层内核。
\end{itemize}

\subsection{应用}

GeLU 是最适合讲 fusion 的逐元素案例，Softmax 是最适合讲行归约的案例，MatMul 是最适合讲 tiling 和 shared memory 的案例。三者分别对应三种常见的 GPU 优化思路。

\begin{table}[H]
\centering
\begin{tabular}{p{3.2cm}p{3.8cm}p{5.4cm}}
\toprule
\textbf{算子} & \textbf{主要瓶颈} & \textbf{课程里的优化手段} \\
\midrule
GeLU & 多次逐元素读写 & fusion, compiler-generated kernel \\
Softmax & 行归约 + 稳定性 & row-wise block, one-pass read/write \\
MatMul & 数据复用和调度 & tiling, shared memory, L2-friendly ordering \\
\bottomrule
\end{tabular}
\caption{三类算子对应三种不同的优化思路}
\end{table}

\subsection{局限}

本讲的方法不是“到处写内核”，而是“先让工具帮你做，再在真正的瓶颈上手写”。很多性能问题会随着库版本、硬件和 shape 变化而变化，所以理论分析只能给方向，不能替代实测。

\begin{knowledgebox}{局限也是方法的一部分}
系统课里最危险的事情不是你不会写内核，而是你以为自己已经知道瓶颈在哪里。这个错误会直接把优化时间浪费掉。
\end{knowledgebox}

%% ========== Section 1: 引言 ==========

